\documentclass[11pt,a4paper]{article}

% =============================================================================
% FEHRADVICE STYLE GUIDE IMPLEMENTATION
% Based on: Appendix SG (REF-STYLE) - Corporate Style Guide
% =============================================================================

% Packages
\usepackage[utf8]{inputenc}
\usepackage[T1]{fontenc}
\usepackage[ngerman]{babel}
\usepackage[
    top=25mm,
    bottom=20mm,
    left=25mm,
    right=20mm
]{geometry}
\usepackage{graphicx}
\usepackage{xcolor}
\usepackage{booktabs}
\usepackage{tabularx}
\usepackage{colortbl}
\usepackage{fancyhdr}
\usepackage{lastpage}
\usepackage{tcolorbox}
\usepackage{enumitem}
\usepackage{hyperref}
\usepackage{helvet}
\usepackage{amsmath}
\usepackage{amssymb}
\renewcommand{\familydefault}{\sfdefault}

% =============================================================================
% FEHRADVICE COLORS (from Corporate-Identity-Guide-FAP-EN-231007)
% =============================================================================

% Primärfarben
\definecolor{fadarkblue}{RGB}{2,64,121}      % #024079 - Headlines, Links, Akzente
\definecolor{falightblue}{RGB}{84,158,222}   % #549EDE - Sekundäre Elemente
\definecolor{fadarkgray}{RGB}{37,33,42}      % #25212A - Fliesstext
\definecolor{falightgray}{RGB}{243,245,247}  % #F3F5F7 - Hintergründe

% Sekundärfarben
\definecolor{falilac}{RGB}{161,160,198}      % #A1A0C6 - Akzente, Diagramme
\definecolor{famint}{RGB}{126,189,172}       % #7EBDAC - Akzente, Diagramme
\definecolor{faocher}{RGB}{222,203,63}       % #DECB3F - Hervorhebungen
\definecolor{faorange}{RGB}{222,157,62}      % #DE9D3E - Warnungen, CTA

% Zusätzliche
\definecolor{faborder}{RGB}{224,224,224}     % Rahmenlinien

% =============================================================================
% HYPERREF SETUP
% =============================================================================

\hypersetup{
    colorlinks=true,
    linkcolor=fadarkblue,
    citecolor=fadarkblue,
    urlcolor=fadarkblue
}

% =============================================================================
% PAGE HEADER & FOOTER (FehrAdvice Style)
% =============================================================================

\pagestyle{fancy}
\fancyhf{}
\fancyhead[L]{\small\color{fadarkgray}\textbf{EBF Management Report}}
\fancyhead[R]{\small\color{fadarkgray}Treppennutzung steigern}
\fancyfoot[L]{\small\color{fadarkblue}\textbf{FehrAdvice \& Partners AG}}
\fancyfoot[R]{\small\color{fadarkgray}Seite \thepage\ von \pageref{LastPage}}
\renewcommand{\headrulewidth}{0.5pt}
\renewcommand{\footrulewidth}{0pt}

% First page style
\fancypagestyle{plain}{
    \fancyhf{}
    \fancyfoot[L]{\small\color{fadarkblue}\textbf{FehrAdvice \& Partners AG}}
    \fancyfoot[R]{\small\color{fadarkgray}Seite \thepage\ von \pageref{LastPage}}
    \renewcommand{\headrulewidth}{0pt}
}

% =============================================================================
% SECTION FORMATTING (FehrAdvice Style)
% =============================================================================

\usepackage{titlesec}

% Section: 22pt Bold, Dunkelblau
\titleformat{\section}
    {\sffamily\bfseries\LARGE\color{fadarkblue}}
    {\thesection}
    {1em}
    {}

% Subsection: 16pt Regular, Dunkelblau
\titleformat{\subsection}
    {\sffamily\large\color{fadarkblue}}
    {\thesubsection}
    {1em}
    {}

% =============================================================================
% TCOLORBOX STYLES (FehrAdvice)
% =============================================================================

\tcbset{
    % Info box (blue)
    fainfobox/.style={
        colback=falightgray,
        colframe=fadarkblue,
        coltitle=white,
        fonttitle=\sffamily\bfseries,
        boxrule=1pt,
        arc=0pt,
        title=#1
    },
    % Warning box (orange)
    fawarningbox/.style={
        colback=faorange!10,
        colframe=faorange,
        coltitle=white,
        fonttitle=\sffamily\bfseries,
        boxrule=1pt,
        arc=0pt,
        title=#1
    },
    % Success box (mint)
    fasuccessbox/.style={
        colback=famint!15,
        colframe=famint!80!black,
        coltitle=white,
        fonttitle=\sffamily\bfseries,
        boxrule=1pt,
        arc=0pt,
        title=#1
    }
}

% =============================================================================
% LIST STYLING
% =============================================================================

\setlist[itemize]{
    label=\textcolor{fadarkblue}{\textbullet},
    leftmargin=*,
    itemsep=4pt
}

\setlist[enumerate]{
    label=\textcolor{fadarkblue}{\arabic*.},
    leftmargin=*,
    itemsep=4pt
}

% =============================================================================
% TABLE STYLING
% =============================================================================

\newcommand{\fahead}[1]{\textcolor{white}{\textbf{#1}}}

% =============================================================================
% DOCUMENT
% =============================================================================

\begin{document}

% Set default text color
\color{fadarkgray}

% =============================================================================
% TITLE PAGE
% =============================================================================

\begin{titlepage}
\begin{center}

\vspace*{1cm}

% Logo placeholder (Schlussstein)
\begin{tikzpicture}
    % Simplified Schlussstein symbol
    \fill[fadarkblue] (0,0) -- (0.8,0) -- (0.6,1.2) -- (0.2,1.2) -- cycle;
    \fill[fadarkblue] (-0.3,0) -- (0.1,0) -- (0.1,0.9) -- (-0.1,0.9) -- cycle;
    \fill[fadarkblue] (0.7,0) -- (1.1,0) -- (0.9,0.9) -- (0.7,0.9) -- cycle;
\end{tikzpicture}

\vspace{0.5cm}

{\Large\color{fadarkblue}\textbf{FEHR ADVICE}}

\vspace{0.2cm}

{\small\color{fadarkgray}Behavioral Economics Consultancy Group}

\vspace{2cm}

% Title
{\Huge\color{fadarkblue}\textbf{Treppe statt Lift}\par}

\vspace{0.8cm}

% Subtitle
{\Large\color{falightblue}Management Report\par}

\vspace{1.5cm}

% Horizontal line
{\color{fadarkblue}\rule{\textwidth}{2pt}}

\vspace{1.5cm}

% Description box
\begin{tcolorbox}[fainfobox={Verhaltensbasierte Intervention zur Gesundheitsförderung}]
Wie Mitarbeitende einer Schweizer Versicherung durch evidenzbasierte Nudges zur Treppennutzung motiviert werden können -- ohne Verbote, ohne Zwang.
\end{tcolorbox}

\vspace{2cm}

% Metadata
\begin{tabular}{ll}
\textbf{Kontext:} & Schweizer Versicherung \\
\textbf{Zielgruppe:} & Mitarbeitende (Bürogebäude) \\
\textbf{Modul:} & BCM2\_HLT\_STAIRS \\
\textbf{Version:} & 1.0 \\
\textbf{Datum:} & Januar 2026 \\
\end{tabular}

\vfill

% Footer
{\small\color{fadarkgray}Evidence-Based Framework for Economic and Social Behavior\par}
\vspace{0.3cm}
{\color{fadarkblue}\textbf{FehrAdvice \& Partners AG}\par}

\end{center}
\end{titlepage}

% =============================================================================
% EXECUTIVE SUMMARY
% =============================================================================

\section*{\color{fadarkblue}Executive Summary}

\begin{tcolorbox}[fainfobox={Kernbotschaft}]
\textbf{Menschen nehmen den Lift nicht aus Faulheit, sondern weil er der \textit{Default} ist.} Die Treppe ist unsichtbar, der Lift omnipräsent. Durch gezielte Choice Architecture -- Sichtbarkeit erhöhen, Defaults ändern, soziale Normen aktivieren -- kann die Treppennutzung um 30--50\% gesteigert werden.
\end{tcolorbox}

\subsection*{Das Problem}
\begin{itemize}
    \item Nur 15--20\% der Mitarbeitenden nutzen regelmässig die Treppe
    \item Durchschnittlicher Büroangestellter: <2'000 Schritte/Tag im Gebäude
    \item Gesundheitskosten durch Bewegungsmangel: CHF 2.4 Mrd./Jahr (Schweiz)
    \item Versicherungen propagieren Gesundheit -- leben sie aber intern?
\end{itemize}

\subsection*{Die Lösung}
\begin{itemize}
    \item \textbf{4 evidenzbasierte Interventionen} aus dem EBF-Toolkit
    \item \textbf{Kein Verbot} -- Lift bleibt verfügbar
    \item \textbf{Keine Kosten} für Infrastruktur (nur Kommunikation)
    \item \textbf{Messbare Wirkung} nach 4 Wochen
\end{itemize}

\subsection*{Erwarteter Impact}
\begin{itemize}
    \item \textbf{+35\%} Treppennutzung (basierend auf Meta-Analyse, Nocon et al. 2010)
    \item \textbf{+1'200 Schritte/Tag} pro aktiven Mitarbeitenden
    \item \textbf{Signalwirkung:} Versicherung lebt, was sie predigt
\end{itemize}

\newpage

% =============================================================================
% SECTION 1: VERHALTENSANALYSE
% =============================================================================

\section{Warum Menschen den Lift nehmen}

Die Entscheidung Treppe vs. Lift ist keine bewusste Kosten-Nutzen-Abwägung. Sie geschieht in \textbf{System 1} -- automatisch, unbewusst, kontextgesteuert.

\subsection{Die vier Barrieren}

\begin{table}[h]
\centering
\rowcolors{2}{white}{falightgray}
\begin{tabular}{>{\columncolor{fadarkblue}}c l p{7cm}}
\rowcolor{fadarkblue}
\fahead{\#} & \fahead{Barriere} & \fahead{Erklärung} \\
\textcolor{white}{1} & \textbf{Unsichtbarkeit} & Treppe ist versteckt, Lift ist zentral \\
\textcolor{white}{2} & \textbf{Default} & Lift ist der «Normalfall» \\
\textcolor{white}{3} & \textbf{Soziale Norm} & Andere nehmen auch den Lift \\
\textcolor{white}{4} & \textbf{Zeitperzeption} & Lift erscheint schneller (ist es meist nicht) \\
\end{tabular}
\caption{Die vier verhaltensökonomischen Barrieren}
\end{table}

\begin{tcolorbox}[fawarningbox={Häufiger Fehler}]
\textbf{Information allein wirkt nicht.} Poster mit «Treppe ist gesund!» ändern nichts, wenn die Barrieren 1--4 bestehen bleiben. Die Intervention muss die \textit{Entscheidungsarchitektur} ändern, nicht das Wissen.
\end{tcolorbox}

\subsection{Der Schweizer Kontext ($\Psi_{CH}$)}

\begin{table}[h]
\centering
\rowcolors{2}{white}{falightgray}
\begin{tabular}{>{\columncolor{fadarkblue}}l l l}
\rowcolor{fadarkblue}
\fahead{Faktor} & \fahead{Wert} & \fahead{Implikation} \\
\textcolor{white}{Gesundheitsaffinität} & Hoch (0.78) & Gesundheitsargumente resonieren \\
\textcolor{white}{Eigenverantwortung} & Hoch (0.82) & Keine Verbote, Wahlfreiheit betonen \\
\textcolor{white}{Soziale Konformität} & Mittel (0.55) & Normen wirken, aber subtil \\
\textcolor{white}{Qualitätsbewusstsein} & Sehr hoch (0.89) & «Richtig machen» als Argument \\
\end{tabular}
\caption{Relevante Kontext-Parameter für die Schweiz (BCM2\_04\_KON)}
\end{table}

\newpage

% =============================================================================
% SECTION 2: INTERVENTIONEN
% =============================================================================

\section{Vier evidenzbasierte Interventionen}

Alle Interventionen sind nach dem \textbf{20-Field Schema} (Appendix HHH) spezifiziert und auf den Schweizer Versicherungskontext kalibriert.

\subsection{Intervention 1: Point-of-Decision Prompts}

\begin{table}[h]
\centering
\rowcolors{2}{white}{falightgray}
\begin{tabular}{>{\columncolor{fadarkblue}}l l}
\rowcolor{fadarkblue}
\fahead{Feld} & \fahead{Spezifikation} \\
\textcolor{white}{F1 Titel} & «Treppenschilder am Entscheidungspunkt» \\
\textcolor{white}{F2 Typ} & Strukturell (Choice Architecture) \\
\textcolor{white}{F3 Kontext} & Organisation (Bürogebäude) \\
\textcolor{white}{F5 FEPSDE} & Physical (P) \\
\textcolor{white}{F6 Journey-Phase} & Action \\
\end{tabular}
\end{table}

\textbf{Umsetzung:} Schilder direkt beim Lift mit Zeitvergleich:

\begin{tcolorbox}[colback=falightgray, colframe=falightblue, boxrule=0.5pt]
\textit{«Treppe: 45 Sekunden bis Stock 3. Lift: 90 Sekunden (Wartezeit + Fahrt).»}
\end{tcolorbox}

\textbf{Evidenz:} Meta-Analyse zeigt +8--12\% Treppennutzung (Nocon et al. 2010)

\subsection{Intervention 2: Soziale Normen sichtbar machen}

\begin{table}[h]
\centering
\rowcolors{2}{white}{falightgray}
\begin{tabular}{>{\columncolor{fadarkblue}}l l}
\rowcolor{fadarkblue}
\fahead{Feld} & \fahead{Spezifikation} \\
\textcolor{white}{F1 Titel} & «Treppennutzungs-Statistik» \\
\textcolor{white}{F2 Typ} & Normativ (Social Proof) \\
\textcolor{white}{F3 Kontext} & Organisation \\
\textcolor{white}{F5 FEPSDE} & Social (S) \\
\textcolor{white}{F6 Journey-Phase} & Awareness $\rightarrow$ Trigger \\
\end{tabular}
\end{table}

\textbf{Umsetzung:} Wöchentliche Anzeige im Eingangsbereich:

\begin{tcolorbox}[colback=falightgray, colframe=falightblue, boxrule=0.5pt]
\textit{«Diese Woche: 847 Kolleg:innen haben die Treppe genommen.»}
\end{tcolorbox}

\textbf{Evidenz:} Soziale Normen erhöhen Treppennutzung um +15--20\% (Burger \& Shelton, 2011)

\subsection{Intervention 3: Ästhetische Aufwertung}

\begin{table}[h]
\centering
\rowcolors{2}{white}{falightgray}
\begin{tabular}{>{\columncolor{fadarkblue}}l l}
\rowcolor{fadarkblue}
\fahead{Feld} & \fahead{Spezifikation} \\
\textcolor{white}{F1 Titel} & «Treppen-Makeover» \\
\textcolor{white}{F2 Typ} & Strukturell (Salienz) \\
\textcolor{white}{F3 Kontext} & Organisation \\
\textcolor{white}{F5 FEPSDE} & Emotional (E) + Physical (P) \\
\textcolor{white}{F6 Journey-Phase} & Maintenance \\
\end{tabular}
\end{table}

\textbf{Umsetzung:}
\begin{itemize}
    \item Bessere Beleuchtung im Treppenhaus
    \item Motivierende Sprüche auf den Stufen («Noch 12... 11... 10...»)
    \item Kunst oder Fotos an den Wänden
\end{itemize}

\textbf{Evidenz:} Ästhetische Aufwertung erhöht Nutzung um +10--15\% (Boutelle et al. 2001)

\newpage

\subsection{Intervention 4: Gamification mit Team-Challenge}

\begin{table}[h]
\centering
\rowcolors{2}{white}{falightgray}
\begin{tabular}{>{\columncolor{fadarkblue}}l l}
\rowcolor{fadarkblue}
\fahead{Feld} & \fahead{Spezifikation} \\
\textcolor{white}{F1 Titel} & «Stockwerk-Challenge» \\
\textcolor{white}{F2 Typ} & Hybrid (Normativ + Feedback) \\
\textcolor{white}{F3 Kontext} & Organisation \\
\textcolor{white}{F5 FEPSDE} & Social (S) + Physical (P) \\
\textcolor{white}{F6 Journey-Phase} & Trigger $\rightarrow$ Maintenance \\
\end{tabular}
\end{table}

\textbf{Umsetzung:} 4-wöchige Abteilungs-Challenge:
\begin{itemize}
    \item Teams sammeln «Stockwerke» (via App oder Strichliste)
    \item Wöchentliches Ranking
    \item Gewinnerteam: Symbolischer Preis (z.B. Team-Lunch)
\end{itemize}

\begin{tcolorbox}[fawarningbox={Crowding-Out Risiko}]
\textbf{Keine monetären Anreize!} Geldprämien würden die intrinsische Motivation (Gesundheit, Team) verdrängen ($\gamma_{\text{Social,Financial}} = -0.2$). Symbolische Anerkennung verstärkt die soziale Norm.
\end{tcolorbox}

% =============================================================================
% SECTION 3: IMPLEMENTIERUNG
% =============================================================================

\section{Implementierungsplan}

\subsection{Portfolio-Logik}

Die vier Interventionen sind \textbf{komplementär} ($\gamma > 0$) und decken verschiedene Journey-Phasen ab:

\begin{table}[h]
\centering
\rowcolors{2}{white}{falightgray}
\begin{tabular}{>{\columncolor{fadarkblue}}l l l}
\rowcolor{fadarkblue}
\fahead{Intervention} & \fahead{Phase} & \fahead{Komplementarität} \\
\textcolor{white}{1. Point-of-Decision} & Action & -- \\
\textcolor{white}{2. Soziale Normen} & Awareness/Trigger & $\gamma_{1,2} = +0.3$ \\
\textcolor{white}{3. Ästhetik} & Maintenance & $\gamma_{1,3} = +0.2$ \\
\textcolor{white}{4. Team-Challenge} & Trigger/Maintenance & $\gamma_{2,4} = +0.4$ \\
\end{tabular}
\caption{Komplementarität der Interventionen}
\end{table}

\subsection{Zeitplan}

\begin{table}[h]
\centering
\rowcolors{2}{white}{falightgray}
\begin{tabular}{>{\columncolor{fadarkblue}}l l l}
\rowcolor{fadarkblue}
\fahead{Woche} & \fahead{Aktion} & \fahead{Verantwortlich} \\
\textcolor{white}{0} & Baseline messen (Zähler installieren) & Facility Management \\
\textcolor{white}{1} & Schilder aufhängen (Intervention 1) & Kommunikation \\
\textcolor{white}{2} & Statistik-Display aktivieren (Int. 2) & HR/Kommunikation \\
\textcolor{white}{3} & Treppen-Makeover (Intervention 3) & Facility Management \\
\textcolor{white}{4--8} & Team-Challenge starten (Int. 4) & HR \\
\textcolor{white}{9} & Evaluation \& Anpassung & Projektleitung \\
\end{tabular}
\caption{Implementierungszeitplan}
\end{table}

\subsection{Erfolgsmessung}

\begin{itemize}
    \item \textbf{Primär:} Treppennutzung (automatische Zähler)
    \item \textbf{Sekundär:} Mitarbeiter:innenzufriedenheit (Pulse Survey)
    \item \textbf{Langfristig:} Krankheitstage, Gesundheitskosten
\end{itemize}

\begin{tcolorbox}[fasuccessbox={Erwartetes Ergebnis}]
\textbf{Prognose:} +35\% Treppennutzung nach 8 Wochen (95\% CI: [+25\%, +45\%])

Basierend auf: Meta-Analyse von 25 Studien (Nocon et al. 2010), adjustiert für Schweizer Kontext ($\Psi_{CH}$) und Versicherungsbranche.
\end{tcolorbox}

\newpage

% =============================================================================
% QUICK REFERENCE CARD
% =============================================================================

\section*{\color{fadarkblue}Quick Reference Card}

\begin{tcolorbox}[fainfobox={Checkliste für das Projektteam}]
\textbf{Vor dem Start:}
\begin{itemize}
    \item[$\square$] Baseline-Messung durchgeführt
    \item[$\square$] Facility Management informiert
    \item[$\square$] Budget für Schilder/Ästhetik genehmigt
    \item[$\square$] Challenge-App/Tool bereit
\end{itemize}

\textbf{Während der Umsetzung:}
\begin{itemize}
    \item[$\square$] Wöchentliche Datenerhebung
    \item[$\square$] Statistik-Display aktuell halten
    \item[$\square$] Challenge-Ranking kommunizieren
\end{itemize}

\textbf{Nach 8 Wochen:}
\begin{itemize}
    \item[$\square$] Evaluation durchführen
    \item[$\square$] Learnings dokumentieren
    \item[$\square$] Entscheidung: Fortführung/Anpassung
\end{itemize}
\end{tcolorbox}

\vspace{1cm}

\begin{tcolorbox}[fainfobox={Die wichtigsten Prinzipien}]
\begin{enumerate}
    \item \textbf{Default ändern, nicht informieren} -- Der Lift ist das Problem, nicht das Wissen
    \item \textbf{Sichtbarkeit vor Motivation} -- Was man nicht sieht, wählt man nicht
    \item \textbf{Sozial vor monetär} -- Anerkennung wirkt stärker als Geld
    \item \textbf{Messen vor Meinen} -- Baseline $\rightarrow$ Intervention $\rightarrow$ Vergleich
\end{enumerate}
\end{tcolorbox}

\vfill

\begin{center}
\textit{Basierend auf EBF Appendix HHH (METHOD-TOOLKIT) und BCM2\_04\_KON (Schweizer Kontext)}

\vspace{0.5cm}

{\color{fadarkblue}\rule{0.5\textwidth}{1pt}}

\vspace{0.5cm}

\textcopyright{} 2026 \textbf{\color{fadarkblue}FehrAdvice \& Partners AG}

\vspace{0.3cm}

{\small\color{fadarkgray}Behavioral Economics Consultancy Group}
\end{center}

\end{document}
